%
%
%
%
%
%
%
\section{Small Bases and Coverage by Dimension}\label{sec:frames}
The constructions of Sections~\ref{sec:residue}--\ref{sec:antipode} need room in the base. A residue sequence needs a check vector with $\ljit(\Lambda_v)>0$, which exists only for $p>2^s$ (Corollary~\ref{cor:residueQU}); the staggered residue sequence has a positive bound only for $p>2^s$ (Corollary~\ref{cor:dyadic}); the block-periodic sequences start at $p=3$, $5$, $7$, $19$ for $s=2,3,4,5$ (Theorem~\ref{thm:block}); window constructions need several residues per window. In the bases $2$, $3$ and $4$ none of them gives anything at every full level: base $2$ is too small for all of them, and in base $3$ the brick-wall bound $(p-3)/(2p)$ is $0$.

This section takes a different route for the smallest bases. Instead of choosing responses, it chooses frames and multipliers for which the map $l\mapsto D_i(l)$ is computed by a small finite automaton reading the base-$q$ digits of $l$. The set of all responses is then finite and known exactly, and the digit relaxation at every level is one finite minimum. This gives closed forms in base~$2$ in dimension two (Theorem~\ref{thm:B2}, which answers a question of Dick, Goda and Suzuki), base~$3$ in dimensions two and three, and base~$4$ in dimensions three and four. The price is that nothing is chosen digit by digit, and the constants fall quickly with the dimension. Section~\ref{sec:primepower} reads the staggered residue sequence over prime-power fields, and Section~\ref{sec:coverage} collects, dimension by dimension, which bases are covered and for which level family.

\subsection{Lacunary Roots and Exact Response Sets}\label{sec:automata}
The responses are coefficients of power series built from the local roots (Proposition~\ref{prop:law}), so their structure is that of the roots (Figure~\ref{fig:automaton}(a)). In the Faure frame of an odd base the root $1+u$, $u+u^2=x$, has the signed Catalan numbers as coefficients, with no visible pattern. In characteristic two the same equation $u^2+u=x$ is solved by the lacunary series $u=\sum_kx^{2^k}$, whose coefficients can be read off the binary digits of the exponent. Every algebraic power series over $\F_q$ is $q$-automatic by Christol's theorem \cite{Chr79,CKMR80}, so automata exist for every algebraic scramble; what a usable bound needs is that they are explicit and small. We use two families of frames whose roots are lacunary. In both, the places form one orbit of a group of affine maps $z\mapsto\alpha z+\beta$ that fix $w$ and $\infty$, so every place looks like every other.

\emph{Artin--Schreier frames.} Let $H\subseteq\F_q$ be an additive subgroup with $|H|=s$ and let the places be the elements of~$H$. Then $w=\wp_H(z)=\prod_{h\in H}(z-h)=\sum_i\wp_{H,i}\,z^{p^i}$ is an additive polynomial, with $\wp_{H,0}\ne0$, and the translations by $H$ fix it. For $H=\F_p$ these are the places $0,\dots,p-1$ of the Faure sequence in dimension $s=p$.

\emph{Kummer frames.} Let $s=q-1$ and let the places be the unit group, $a_c=\zeta^c$ for a generator $\zeta$ of $\F_q^*$, with $w=z^{q-1}-1$; the multiplications by $\zeta$ fix it. For $q=3$ the places are $\pm1$, and $y=2z-1$ carries the Faure places $\{0,1\}$ onto these ($y^2-1=4(z^2-z)$ in $\F_3$).

\begin{lemma}[Lacunary roots]\label{lem:frames}
Let $q$ be a power of the prime $p$.
\begin{enumerate}
\item[(a)] In the Artin--Schreier frame of $H$ the local roots are $z_c=a_c+U$, where $U\in x\F_q[[x]]$ is the solution of $\wp_H(U)=x$. It is supported on the powers of $p$ and satisfies $U(x)=U(x^q)-\wp^\flat_H(x)$, where $\wp^\flat_H$ is the additive polynomial with $\wp^\flat_H(\wp_H(z))=z^q-z$. For $H=\F_q$, $U=-\sum_{k\ge0}x^{q^k}$.
\item[(b)] In the Kummer frame the local roots are $z_c=a_c\tau$, where $\tau^{q-1}=1+x$, $\tau(0)=1$, and
\[
\tau^{-1}=\prod_{i\ge0}\bigl(1+x^{q^i}\bigr),\qquad \tau^{-1}(x)=(1+x)\,\tau^{-1}(x^q);
\]
$\coef{n}\tau^{-1}=1$ if every base-$q$ digit of $n$ is $0$ or $1$, and $0$ otherwise.
\end{enumerate}
In both cases the root has degree $s$ over $\F_q(x)$.
\end{lemma}

\begin{proof}
(a) $\wp_H(a_c+U)=\wp_H(U)=x$, so $a_c+U$ is the local root at $a_c$. The recursion $U=\wp_{H,0}^{-1}\bigl(x-\sum_{i\ge1}\wp_{H,i}U^{p^i}\bigr)$ preserves support in the powers of $p$. Since $\wp^\flat_H\circ\wp_H$ is monic of degree $q$ and vanishes on $\F_q$, it is $z^q-z$; hence $U^q-U=\wp^\flat_H(x)$, and $U^q=U(x^q)$. (b) $T=\prod_i(1+x^{q^i})$ satisfies $T(x)^q=T(x^q)=T(x)/(1+x)$, so $\tau=T^{-1}$ solves $\tau^{q-1}=1+x$, and $(a_c\tau)^{q-1}=\tau^{q-1}$. Expanding the product, $x^n$ occurs once for each way of writing $n$ as a sum of distinct powers of $q$. The degree: $\wp_H(Z)-x$ is linear in $x$, and $Z^{q-1}-(1+x)$ is Eisenstein at $1+x$.
\end{proof}

Both relations are \emph{Mahler equations}: the root at $x$ is an affine function of the root at $x^q$, additive for translations and multiplicative for rotations. We take $\theta=1$ and multipliers that are rational functions of the root, $g_c\in\F_q(x)(U)$ resp.\ $\F_q(x)(\tau)$. A scramble of these frames is \emph{symmetric} if all coordinates use the same multiplier $g_c=g$; then $V(z)\mapsto V(z+h)$ (resp.\ $V(\zeta z)$) permutes the coordinates of every point and preserves every prefix.

To read a coefficient digit by digit we use the \emph{Cartier operators} $\mathcal C_r\bigl(\sum_nc_nx^n\bigr)=\sum_nc_{qn+r}x^n$, $0\le r<q$. They satisfy $\mathcal C_r\bigl(A(x^q)F\bigr)=A\,\mathcal C_rF$ and
\begin{equation}\label{eq:cartierread}
\coef{l}F=\coef{0}\,\mathcal C_{l_k}\cdots\mathcal C_{l_1}\mathcal C_{l_0}F\qquad\text{for }l=\textstyle\sum_{i\le k}l_iq^i.
\end{equation}
For example, $q=3$ and $F=\tau^{-1}$: by Lemma~\ref{lem:frames}(b), $\mathcal C_r\tau^{-1}=\mathcal C_r(1+x)\cdot\tau^{-1}=[r\le1]\,\tau^{-1}$, so reading the digits of $n$ visits two states, $\tau^{-1}$ and $0$, and a digit $2$ kills the coefficient (Figure~\ref{fig:automaton}(c)). The responses behave in the same way. By Proposition~\ref{prop:law}, $D_i(l)=\coef{l}\mathcal H_i$ with
\[
\mathcal H_i:=x^{-i}\FM_{>i}\hat G=\frac{x^{-i}\bigl(I-\FM_{\le i}G\bigr)}{1-x},\qquad \FM_{\le i}=\textstyle\sum_{a\le i}\FM_ax^a .
\]

\begin{proposition}[Exact response sets]\label{prop:automaton}
Let a self-similar scramble have a frame matrix $\FM(x)$ with entries in $\F_q(x)(\vartheta)$, where $\vartheta\in\F_q[[x]]$ has degree $s'$ over $\F_q(x)$ and $\vartheta(x)=A(x)\,\vartheta(x^q)+B(x)$ with $A,B\in\F_q[x]$. Fix $j\ge1$ and let $\vec{\mathcal H}$ be the vector of the series $\mathcal H_i$ and $x\mathcal H_i$, $i<j$. There is a finite set $\mathcal S$ of such vectors, containing $\vec{\mathcal H}$ and stable under every $\mathcal C_r$, with
\[
\bigl(D_i(l),\,D_i(l-1)\bigr)_{i<j}=\coef{0}\,\mathcal C_{l_k}\cdots\mathcal C_{l_0}\vec{\mathcal H}\qquad(l\ge1).
\]
So the set $\{(D_i(l),D_i(l-1))_{i<j}:l\ge1\}$ is exactly the set of values $\coef{0}S$ over the states $S\in\mathcal S$ reached from $\vec{\mathcal H}$ by a digit word with a nonzero digit, and a finite breadth-first search computes it.
\end{proposition}

The proof (Appendix~\ref{app:smallbases}) writes every series over one common denominator and shows that the Cartier operators keep the degrees of the numerators bounded.

\begin{figure}[t]
\centering
\paperfig{fig_automaton_v4}
\caption{Finite carry patterns. (a)~Coefficients $\coef{n}$, $n\le80$, of three local roots (blank $=0$): the Faure root $u$, $u+u^2=x$, over $\F_7$ (signed Catalan numbers modulo $7$), the base-$2$ root $u=\sum_kx^{2^k}$, and the Kummer root $\tau^{-1}$ over $\F_3$, whose coefficient is $1$ exactly when $n$ has no base-$3$ digit $2$. All three are automatic by Christol's theorem; only the lacunary ones have a digit rule small enough for an explicit bound. (b)~Long carries repeat short ones for the base-$2$ sequence (Lemma~\ref{lem:octave}, $j=12$): the responses depend only on the parity of $l$ and the bits $\lfloor\log_2n\rfloor\bmod2$ on $[l,l+12]$. A range crossing a power of two, $[250,262]\ni256$, has the partner $l-3\cdot2^{k-1}=58$, which crosses $64$ at the same relative position; a range inside one octave, $[150,162]$, has a partner $[40,52]$ in an octave of the same parity. (c)~The two-state automaton that reads the base-$3$ digits of $n$ and outputs $\coef{n}\tau^{-1}$ over $\F_3$.}
\label{fig:automaton}
\end{figure}

The proposition applies to the Artin--Schreier frames with $\vartheta=U$, $A=1$, $B=-\wp^\flat_H$, and to the Kummer frames with $\vartheta=\tau^{-1}$, $A=1+x$, $B=0$ (Lemma~\ref{lem:frames}). The pairs $(D_i(l),D_i(l-1))$ are what Lemma~\ref{lem:cert} needs for $q=4$, through (\ref{eq:carrychange}). With all $l\ge1$ in place of $l\le K$, Lemma~\ref{lem:cert} then gives one bound for every level with $q^K\ge4$. The point of the proposition is exactness, not finiteness (which Christol's theorem gives): a screen over carry lengths up to some bound can miss binding responses. For the base-$3$ sequence of Remark~\ref{rem:B3}, the last of its $84$ response tuples first appears at $l=2469$.

\subsection{Dimension Two}\label{sec:smalltwo}
In base two every digital $(0,2)$-sequence is, up to a reordering inside each prefix, a scramble $(L_0,L_1\Pas_1)$ of the Sobol' sequence \cite{HS19}. Dick, Goda and Suzuki prove that neither one-sided choice, $(L,\Pas_1)$ or $(\Pas_0,L\Pas_1)$, is well separated, and ask whether a two-sided pair is \cite[Section~4.2]{DGS2}; random linear scrambles of binary digital nets with fixed $t$ have mesh ratio $O_P(\log N)$ under a uniform coincidence condition, which the Sobol' prefixes satisfy \cite[Thm.~5.5, Lem.~4.10]{Suz26}. The one-sided failure is an instance of Corollary~\ref{cor:unscrambled}: a coordinate that is left a radical inverse up to a scalar forces close pairs. The following explicit pair, which scrambles both coordinates, answers the question positively.

\begin{theorem}[Base 2]\label{thm:B2}
Let $u=\sum_{k\ge0}x^{2^k}\in\F_2[[x]]$, so that $u^2+u=x$, and let the digital sequence over $\F_2$ have the digits
\[
y_{0,n}(V)=\coef{n}\,(1+x)(1+u)\,V(u),\qquad y_{1,n}(V)=\coef{n}\,\frac ux\,V(1+u),
\]
that is, $C_0=L_0$ and $C_1=L_1\Pas_1$ with $L_0[n][k]=\coef{n}(1+x)(1+u)u^k$ and $L_1[n][k]=\coef{n+1}u^{k+1}$. It is a $(0,2)$-sequence in base~$2$, and for every $K\ge1$
\[
\sep{P_{2K}}\ \ge\ \frac{\sqrt{2269618}}{4096}=0.36780\ldots
\]
Consequently $\delta(P_{[N]})\ge0.1839\,N^{-1/2}$ for every $N\ge2$, $h/\delta\le7.7$ along all $N$ ($15.4$ in the convention of \cite{DGS2}), and every prefix $P_m$ is $k_m$-separated in the sense of \cite{DGS2} (their $\kappa_m$) with $k_m=\lceil m/2\rceil+3$. In particular the sequence is quasi-uniform.
\end{theorem}

This is the Faure frame in base $2$ ($w=z^2+z$, roots $z_0=u$, $z_1=1+u$, $\theta=1$) with $g_0=(1+x)(1+u)$ and $g_1=u/x=1/(1+u)$, an algebraic but not symmetric scramble of the Artin--Schreier frame $H=\F_2$. The multipliers are chosen so that one row of the inverse frame matrix is rational and the other lacunary:
\begin{equation}\label{eq:B2G}
\FM=\begin{pmatrix}g_0&g_0u\\ g_1&g_1(1+u)\end{pmatrix},\qquad G=\FM^{-1}=\begin{pmatrix}\frac1{1+x}&x\\[2pt] \frac u{x(1+x)}&1+u\end{pmatrix}.
\end{equation}
The one-digit responses are $O_0^{(l)}=l\bmod2$, $X_{0\leftarrow1}^{(l)}=1$, $O_1^{(l)}=\lfloor\log_2l\rfloor\bmod2$ and $X_{1\leftarrow0}^{(l)}=\sum_{n<l}\lfloor\log_2(n+2)\rfloor\bmod2$. In base two the automaton of Proposition~\ref{prop:automaton} can be replaced by a hand argument: long carries repeat the patterns of short ones.

\begin{lemma}[Long carries repeat short ones]\label{lem:octave}
For the sequence of Theorem~\ref{thm:B2} and $j\le12$, the tuples $(D_i(l))_{i<j}$, $l\ge1$, take the same values as for $1\le l\le63$; for $j=12$ there are $35$ of them.
\end{lemma}

The proof (Appendix~\ref{app:smallbases}) reads the entries of $\hat G$ off the binary digits and maps every carry length $l\ge64$ to a smaller one with the same responses.

\begin{proof}[Proof of Theorem~\ref{thm:B2}]
The sequence is a self-similar scramble of the Faure frame, so it is a $(0,2)$-sequence (Lemma~\ref{lem:t0}) whose offset map is the same at every level (Proposition~\ref{prop:law}). By Lemma~\ref{lem:octave}, Lemma~\ref{lem:cert} with $j=12$ is a minimum over $35+35+2\cdot35^2$ cases (axis steps with one carry length, diagonal and antidiagonal steps with two). This exact computation (computer-assisted) gives $\LB_{12}=\sqrt{2269618}/4096$, attained at the horizontal step with $l=1$: the own change $\sigma_0=111110011000$ and the cross change $\sigma_1=101000100000$ give $f_{12}(1,\sigma_0)=103/4096$ and $f_{12}(0,\sigma_1)=1503/4096$. This covers $K\ge2$; for $K=1$, $\sep{P_2}=0.6334$. Lemma~\ref{lem:prefix} with $q=2$ and Remark~\ref{rem:prefixcases} give $\delta(P_{[N]})\ge(\LB_{12}/2)N^{-1/2}$ and $h/\delta\le2\sqrt2/\LB_{12}<7.7$. Finally two points in a square of side $2^{-k_m}$ are less than $\sqrt2\cdot2^{-k_m}$ apart, and $\sqrt2\cdot2^{-\lceil m/2\rceil-3}<0.1839\cdot2^{-m/2}\le\delta(P_m)$.
\end{proof}

Since $k_m\le m/2+4$ and $\delta$ is measured on the torus, the sequence is well-separated, and toroidally well-separated, in the sense of \cite[Def.~3.8]{DGS2}, which is the property asked for in \cite[Section~4.2]{DGS2}.

\begin{figure}[t]
\centering
\paperfig{fig_B2_v4}
\caption{XOR staggering in base~$2$ (Theorem~\ref{thm:B2}). (a)~$12\times12$ of the $64\times64$ cells of $P_{12}$, with lattice images and jitter; cell boundaries shaded by carry length. The closest pair is a horizontal neighbour pair across a boundary of carry length $1$, staggered vertically by the cross mask. (b)~The masks $D_i(l)$ XORed onto the first six geometric digits (rows, worth $2^{-1},\dots,2^{-6}$ of a cell) for $l\le64$; the strips on top give the pattern of $l$ (Lemma~\ref{lem:octave}). Boxed: the binding case, a horizontal step with $l=1$, with the $12$-digit masks $\sigma_0=111110011000$ (own) and $\sigma_1=101000100000$ (cross). (c)~The guaranteed distance across a boundary of carry length $l$ (Lemma~\ref{lem:cert}, $j=12$); the minimum $0.3678$ is the horizontal step with $l=1$. Exact $\sep{P_{12}}=0.3968$.}
\label{fig:B2}
\end{figure}

\emph{Why it works.} The offset map of a base-two sequence does not shift points by small real amounts: what is fixed per carry length is an $\F_2$-mask XORed onto the geometric digits (Figure~\ref{fig:B2}(b)). At the most common boundary, $l=1$ in coordinate~$0$, the own mask can pull the neighbour almost onto the vertical line through the point. But the cross mask has a leading~$1$: the vertical offsets differ in their first digit, by $\pm\frac12$ of a cell, corrected by the lower digits by at most $\frac18+\frac1{128}$ plus the jitter, so the two points stay at least $1503/4096>0.366$ of a cell apart vertically. Likewise every step of coordinate~$1$ flips the leading horizontal digit ($X_{0\leftarrow1}^{(l)}=1$). This is the brick wall of Example~\ref{ex:brickwall} in XOR form: where the own mask could close the gap, the cross mask forces a stagger of about half a cell. The true separation is somewhat larger: $\sep{P_m}=0.3968$ for every even $m$ from $12$ to $26$, against $0.0442$ for the Sobol' sequence at $m=12$ and $0.0007$ at $m=24$ (exact). The generator data, as a drop-in table of $32$-bit direction numbers with a reference implementation, are in the archive (Appendix~\ref{app:archive}).

\begin{remark}[Base 3]\label{rem:B3}
In base three the Faure frame is the Kummer frame $\{\pm1\}$, with the lacunary root $\tau=\sqrt{1+x}$ of Lemma~\ref{lem:frames}(b), and Proposition~\ref{prop:automaton} handles it as well. Take the roots $z_0=\tau-1$, $z_1=2-\tau$ of $z^2-z=x$, $\theta=1$ and the multipliers $g_0=-(1+x)(1-x+\tau)/(1-x)$, $g_1=-(1+\tau)/\tau$. This self-similar scramble $(L_0\Pas_0,L_1\Pas_1)$ of the Faure sequence in base $3$ is a $(0,2)$-sequence (Lemma~\ref{lem:t0}), and Proposition~\ref{prop:automaton} with $\vartheta=\tau^{-1}$, $A=1+x$, $B=0$, $j=6$ (a search of $190$ states) gives exactly $84$ response tuples; Lemma~\ref{lem:cert} over all steps and pairs of tuples gives $\LB_6=0.266926$, attained at an antidiagonal step (computer-assisted). With $\sep{P_2}=0.5828$ this gives $\sep{P_{2K}}\ge0.2669$ for every $K\ge1$: base three in dimension two at every full level, where Theorem~\ref{thm:block}(a) gives $\sep{P_{4K}}\ge1/3$ on the sparser level family $4\N$. The unscrambled Faure sequence in base~$3$ is not quasi-uniform \cite{GHS26} (Corollary~\ref{cor:faure}); its separation falls to $0.016$ at $m=9$, while that of this scramble stays at $0.298$ or more for even $m\le14$ (exact).
\end{remark}

\begin{corollary}[Every prime base, dimension two]\label{cor:allprimes}
For every prime $p$ there is an explicit closed-form scramble $(L_0\Pas_0,L_1\Pas_1)$ of the Faure sequence in base $p$ that is a quasi-uniform $(0,2)$-sequence with a bound at every full level: Theorem~\ref{thm:B2} for $p=2$, Remark~\ref{rem:B3} for $p=3$ and the brick wall (Theorem~\ref{thm:WF}) for $p\ge5$, with $\sep{P_{2K}}\ge0.3678$, $0.2669$ and $(p-3)/(2p)$, respectively, for every $K$.
\end{corollary}

For large primes the hexagonal brick wall improves the constant at every full level to $1-17/p$ (Corollary~\ref{cor:HB17}, $p>3677$), and for every prime $p\ge3$ the block-periodic residue sequence reaches $1-2/p$ at every level $m\in4\N$ (Theorem~\ref{thm:block}(a)).

\subsection{Dimensions Three and Four, Small Bases}\label{sec:smallthree}
The same method works with symmetric scrambles of the Artin--Schreier frames $H=\F_3$ and $H=\F_4$ and of the Kummer frame $\F_4^*$. In each case the multiplier is a short rational function of the root, found by an exhaustive screen over small families of such functions; the proofs are identical, and their finite parts are collected in Table~\ref{tab:frames}.

\begin{theorem}[Base 3, dimension 3]\label{thm:S3}
Let $U=-\sum_{k\ge0}x^{3^k}\in\F_3[[x]]$, so that $U^3-U=x$, and $g=(1+x-U)/(1-x)$. The digital sequence over $\F_3$ with the digits $y_{c,n}(V)=\coef{n}\,g(x)\,V\bigl(c+U(x)\bigr)$, $c\in\F_3$, is a scramble of the Faure sequence $(\Pas_0,\Pas_1,\Pas_2)$ in base~$3$ and a $(0,3)$-sequence, and $\sep{P_{3K}}\ge0.05598$ for every $K\ge1$. Consequently $\delta(P_{[N]})\ge0.0186\,N^{-1/3}$ for all $N\ge2$ and $h/\delta<93$.
\end{theorem}

\begin{theorem}[Base 4, dimension 3]\label{thm:K4}
Let $\F_4=\{0,1,\omega,\omega^2\}$, let $\tau^{-1}=\prod_{i\ge0}(1+x^{4^i})\in\F_4[[x]]$, so that $\tau^3=1+x$, and
\[
g=\frac{\omega^2(\tau+\omega)}{\tau+\omega^2}=1+x+\omega x^2+\omega^2x^3+\omega^2x^5+\cdots.
\]
The digital sequence over $\F_4$ with the digits $y_{c,n}(V)=\coef{n}\,g(x)\,V(\omega^c\tau(x))$, $c=0,1,2$, is a $(0,3)$-sequence in base~$4$, and $\sep{P_{3K}}\ge0.11293$ for every $K\ge1$. Consequently $\delta(P_{[N]})\ge0.0282\,N^{-1/3}$ for all $N\ge2$ and $h/\delta<61.4$. With the binary code of $\F_4$ it is a binary digital sequence, generated by XOR like the Sobol' sequence.
\end{theorem}

\begin{theorem}[Base 4, dimension 4]\label{thm:S4}
Let $U=\sum_{k\ge0}x^{4^k}\in\F_4[[x]]$, so that $U^4+U=x$, and $g=1+\omega U+U^2$. The digital sequence over $\F_4$ with the digits $y_{c,n}(V)=\coef{n}\,g(x)\,V(c+U(x))$, $c\in\F_4$, is a $(0,4)$-sequence in base~$4$ and a binary digital sequence, and $\sep{P_{4K}}\ge0.03013$ for every $K\ge1$. Consequently $\delta(P_{[N]})\ge0.0075\,N^{-1/4}$ for all $N\ge2$ and $h/\delta<266$.
\end{theorem}

\begin{proof}[Proof of Theorems~\ref{thm:S3}, \ref{thm:K4} and~\ref{thm:S4}]
By Lemma~\ref{lem:frames} the frames are the Artin--Schreier frames $H=\F_3$ ($w=z^3-z$) and $H=\F_4$ ($w=z^4+z$) with roots $c+U$, and the Kummer frame $\F_4^*$ ($w=z^3-1$) with roots $\omega^c\tau$; the multiplier is the same rational function of the root in every coordinate. (For Theorem~\ref{thm:K4}, coordinate $c$ reads $V(z)\,\omega^2(z-a_{c+1})/(z-a_{c+2})$ at $a_c=\omega^c$, indices modulo $3$: the multiplier has its zero and its pole at the two other places.) So $t=0$ by Lemma~\ref{lem:t0}, and Proposition~\ref{prop:automaton} gives the exact response sets of Table~\ref{tab:frames}, with the carry correction (\ref{eq:carrychange}) for $q=4$. Lemma~\ref{lem:cert} over all steps and all combinations of response columns gives the values $\LB_j$ of the table (computer-assisted). For $q=4$ this covers every $K\ge1$; for $q=3$ it covers $K\ge2$, and $\sep{P_3}=0.4887$. Lemma~\ref{lem:prefix} gives the rest. With $\F_4=\F_2\oplus\F_2\omega$, multiplication by an element of $\F_4$ is an $\F_2$-linear $2\times2$ block and addition is XOR of the codes, which gives the binary form.
\end{proof}

\begin{table}[t]
\centering
\small
\caption{The finite parts of the small-base theorems: depth $j$, number of distinct responses (per moving coordinate; in brackets the change columns after the $\F_4$ carry correction), bound $\LB_j$ for every level, the step attaining it (up to the symmetry of the frame), and the exact separation at the largest level we computed. All entries are exact computations; the responses are exact sets by Proposition~\ref{prop:automaton} (Lemma~\ref{lem:octave} for base~$2$). Level family: all full levels, $\Lev=s\N$.}
\label{tab:frames}
\begin{tabular}{@{}lllclll@{}}
\toprule
result & base, $s$ & frame & $j$ & responses & $\LB_j$ & minimising step; exact $\sep{P_m}$\\
\midrule
\ref{thm:B2} & $2$, $2$ & Faure ($H=\F_2$) & $12$ & $35$ tuples & $0.367804$ & $(1,0)$, $l=1$; $0.3968$ ($m\le26$)\\
\ref{rem:B3} & $3$, $2$ & Faure ($\F_3^*$) & $6$ & $84$ tuples & $0.266926$ & antidiagonal; $0.298$ ($m=14$)\\
\ref{thm:S3} & $3$, $3$ & $H=\F_3$ & $8$ & $85$ & $0.055983$ & $(1,-1,-1)$; $0.129$ ($m=15$)\\
\ref{thm:K4} & $4$, $3$ & $\F_4^*$ & $8$ & $32$ ($67$) & $0.112930$ & $(0,1,1)$; $0.168$ ($m=12$)\\
\ref{thm:S4} & $4$, $4$ & $H=\F_4$ & $5$ & $64$ ($120$) & $0.030131$ & all four move; $0.101$ ($m=12$)\\
\bottomrule
\end{tabular}
\end{table}

Dick, Goda and Suzuki prove that the Faure sequence in base~$3$ is not quasi-uniform \cite{DGS2}, and Goda, Hofer and Suzuki extend this to scrambles $(\Pas_0,L_1\Pas_1,L_2\Pas_2)$ that leave the first coordinate unscrambled, and likewise in every prime base $p$ in dimension $p$ \cite[Rem.~3]{GHS26} (Corollary~\ref{cor:unscrambled} contains both). Theorem~\ref{thm:S3} scrambles all three coordinates ($g$ is not constant), so it is outside the scope of these results and consistent with them.

The certified bounds in dimensions three and four are far below the exact separations of Table~\ref{tab:frames} because the binding steps involve long carries that the computable prefixes do not contain; explicit close pairs of points bracket the true infimum closely (Appendix~\ref{app:sharp}, which also explains why these bases work, Remark~\ref{rem:height}).

\emph{Base seven.} At $p=7$ in dimension three the residue sequences give nothing ($7<2^3$), and the block-periodic residue sequence covers only the levels $m\in9\N$ (Theorem~\ref{thm:block}(b), $1-2/7$). Every full level is covered by a bound of a different kind, in which the admissible own response depends on the current cross responses of its column; it is given in Appendix~\ref{app:p7} (Theorem~\ref{thm:p7}: $\sep{P_{3K}}\ge1/7$ for every $K$, computer-assisted).

\subsection{Prime-Power Bases}\label{sec:primepower}
Apart from $q=4$, the constructions above live over prime fields. The staggered residue sequence does not need this: its two lemmas use only the residue theorem, which holds over every field. What changes over $\F_q$, $q=p^r$ with $r\ge2$, is the carry. Adding $1$ to the integer cell index changes the finest digit by an element $\Delta(y)$ that depends on the digit (Remark~\ref{rem:F4carry}), not by the constant~$1$, so the rank-1 rule of the staggered coordinate is read with several possible increments. The resulting bound is a small finite minimum, and it depends on how the digits of $\F_q$ are encoded as integers.

Write $\F_q=\F_p[x]/(f)$ for a monic irreducible $f$ of degree $r$, and encode $\sum_{i<r}c_ix^i$ ($0\le c_i<p$) as the integer $\sum_ic_ip^i$ (the power-basis encoding of $f$, a natural identification in the sense of Section~\ref{sec:prelim}). With this encoding $\Delta(y)$ is one of the $r$ elements $1$, $1+x$, \dots, $1+x+\dots+x^{r-1}$, according to the number of trailing digits $p-1$ of the code of $y$; let $\mathcal D_f$ be this set. For weights $k\in(\F_q^*)^{s-1}$ put
\[
\beta_q(k)=\min\ f_1\Bigl(e_{\cstar},\ \sum_{c\in C}e_c\,k_c\,\Delta_c\Bigr),
\]
the minimum over all steps $e\in\{-1,0,1\}^s\setminus\{0\}$ and all $\Delta_c\in\mathcal D_f$ for the moving coordinates $c\in C$, with $f_1$ the one-digit function of Lemma~\ref{lem:cert} for this encoding.

\begin{proposition}[Prime-power bases]\label{prop:primepower}
Let $(s,q,f,k)$ be a row of Table~\ref{tab:primepower}. Read Definition~\ref{def:residuestar} over $\F_q$ (any $s$ distinct places in $\F_q$, any $\theta,\mu_\star\in\F_q^*$, the weights $k$), with the power-basis encoding of $f$. The result is a digital $(0,s)$-sequence over $\F_q$ with
\[
\sep{P_{sK}}\ \ge\ \beta_q(k)\qquad\text{for every }K\ge1,
\]
with the value of the table; it is quasi-uniform, with $\delta(P_{[N]})\ge(\beta_q(k)/q)\,N^{-1/s}$ for every $N\ge2$. The evaluation of $\beta_q(k)$ is computer-assisted.
\end{proposition}

\begin{proof}
The proofs of Lemmas~\ref{lem:residuestar-t0} and~\ref{lem:residuestar-rel} use Lemma~\ref{lem:resform} and Lemma~\ref{lem:t0}, which hold over every finite field; so the sequence is a digital $(0,s)$-sequence over $\F_q$, and $y_{\cstar,K}=\sum_{c\in C}k_c\,y_{c,K-1}$ on $P_{sK}$ for every $K$. In the language of Section~\ref{sec:law} the responses of the staggered coordinate are constant, $D(l)[\cstar][c]=k_c$ and $D(l)[\cstar][\cstar]=0$ for every $l\ge1$. By (\ref{eq:carrychange}) a step of coordinate $c$ with carry length $l$, stopping at the digit $y$, therefore changes $r_{\cstar}$ by $\Delta(\bar y)k_c$ if $l\ge2$ (the wrap included) and by $\Delta(y)k_c$ if $l=1$; in both cases by $\Delta_ck_c$ with $\Delta_c\in\mathcal D_f$. Lemma~\ref{lem:cert} with $j=1$, keeping only the term of coordinate $\cstar$, gives $\sep{P_{sK}}\ge\min(1,\beta_q(k))=\beta_q(k)$ for every $K$ with $q^K\ge4$, which is every $K\ge1$ since $q\ge8$. The value $\beta_q(k)$ is a minimum over at most $3^s\,r^{s-1}$ terms, each a minimum over the $q$ positions of $\enc$ differences, evaluated exactly (computer-assisted). The last statement is Lemma~\ref{lem:prefix} with $\Lev=s\N$.
\end{proof}

\begin{table}[t]
\centering
\footnotesize
\caption{Proposition~\ref{prop:primepower}: the staggered residue sequence over $\F_q$, $q=p^r$, with the power-basis encoding of the monic irreducible $f$; weights $k$ by their codes; $\beta_q$ is certified for every $K$ ($\Lev=s\N$). For $s\le3$, $f$ and $k$ give the best value over all $f$ and $k$ (exhaustive); for $s=4,5$ they were found by search. Last column: exact $\sep{P_3}$ for the places $0,1,2$ (codes), $\cstar=2$, $\theta=1$. $\beta_q$ depends on the encoding: for $q=25$, $s=2$ it ranges from $5/25$ ($x^2+2$) to $8/25$ ($x^2+3$), and for $q=32$, $s=3$ it is positive for $2$ of the $6$ irreducible quintics.}
\label{tab:primepower}
\setlength{\tabcolsep}{2.5pt}
\begin{tabular}[t]{@{}r l l r@{}}
\toprule
\multicolumn{4}{@{}l}{$s=2$}\\
$q$ & $f$ & $k$ & $\beta_q$\\
\midrule
8 & $x^3+x+1$ & $(2)$ & $1/8$\\
9 & $x^2+1$ & $(3)$ & $1/9$\\
16 & $x^4+x+1$ & $(2)$ & $1/16$\\
25 & $x^2+3$ & $(10)$ & $8/25$\\
27 & $x^3+x^2+2x+1$ & $(15)$ & $5/27$\\
32 & $x^5+x^3+x^2+x+1$ & $(23)$ & $7/32$\\
49 & $x^2+2x+2$ & $(21)$ & $19/49$\\
64 & $x^6+x^4+x^3+x+1$ & $(45)$ & $13/64$\\
81 & $x^4+x^2+x+1$ & $(31)$ & $17/81$\\
121 & $x^2+9$ & $(55)$ & $53/121$\\
125 & $x^3+3x+3$ & $(51)$ & $44/125$\\
128 & $x^7+x^5+x^4+x^3+x^2+x+1$ & $(95)$ & $31/128$\\
\bottomrule
\end{tabular}\hfill
\begin{tabular}[t]{@{}c r l l r c@{}}
\toprule
$s$ & $q$ & $f$ & $k$ & $\beta_q$ & $\sep{P_3}$\\
\midrule
3 & 25 & $x^2+x+2$ & $(9,12)$ & $2/25$ & $0.1961$\\
 & 32 & $x^5+x^3+x^2+x+1$ & $(12,2)$ & $1/32$ & $0.1814$\\
 & 49 & $x^2+1$ & $(21,7)$ & $5/49$ & $0.3565$\\
 & 64 & $x^6+x+1$ & $(42,23)$ & $2/64$ & $0.1781$\\
 & 81 & $x^4+2x^3+x^2+x+2$ & $(3,9)$ & $2/81$ & $0.1848$\\
 & 121 & $x^2+2x+5$ & $(22,54)$ & $20/121$ & $0.2719$\\
 & 125 & $x^3+4x^2+4$ & $(45,66)$ & $14/125$ & $0.2323$\\
 & 128 & as for $s=2$ & $(56,95)$ & $7/128$ & $0.1767$\\
 & 169 & $x^2+4x+9$ & $(51,89)$ & $35/169$ & $0.2320$\\
\addlinespace
4 & 49 & $x^2+1$ & $(7,16,19)$ & $1/49$ & --\\
 & 121 & $x^2+x+8$ & $(21,31,49)$ & $7/121$ & --\\
 & 125 & $x^3+x+1$ & $(2,5,10)$ & $1/125$ & --\\
\addlinespace
5 & 121 & $x^2+1$ & $(13,28,36,47)$ & $1/121$ & --\\
 & 169 & $x^2+2$ & $(17,58,63,83)$ & $2/169$ & --\\
\bottomrule
\end{tabular}
\end{table}

The bound uses one coordinate and one digit, as in Corollary~\ref{cor:dyadic}, and it is small; the exact separations are several times larger (last column of Table~\ref{tab:primepower}; for $s=3$ the smallest $\sep{P_{[N]}}$ over all $N$ up to $10^7$ is between $0.0144$ ($q=32$) and $0.1484$ ($q=49$), above $\beta_q/q$ in every case). The proof uses the encoding only through the set $\mathcal D_f$ of increments $\Delta(y)$, so it holds for every bijection $\enc$ with $\mathcal D_f$ replaced by its set of increments. A zero is a zero of this bound, not a statement about the base: for $q=8,9$ in dimension three $\beta_q=0$ for every digit code and every weight, for $q=16,27$ for every $\F_p$-linear code; in dimension four $\beta_q=0$ for $q=16,25,27,32,64,81$, and in dimension five for $q=49$, for every power-basis encoding and every weight (exact computations). Other constructions and certificates with more digits are untried there.

\subsection{Coverage by Dimension}\label{sec:coverage}
Which bases admit a quasi-uniform digital $(0,s)$-sequence with a certified bound, and on which level family? The answer collects the results of Sections~\ref{sec:residue}--\ref{sec:frames}. The level family matters: every family of this paper gives quasi-uniformity, but a denser family gives the better bound for every $N$ (Lemma~\ref{lem:prefix}), and the full levels $\Lev=s\N$ are the classical choice. Table~\ref{tab:coverage} and Figure~\ref{fig:dim3} give the details.

\begin{theorem}[Coverage]\label{thm:coverage}
In each of the following cases there is an explicit digital $(0,s)$-sequence over $\F_q$ with a positive lower bound on $\sep{P_m}$ for every $m$ in the stated level family $\Lev$, and hence, by Lemma~\ref{lem:prefix}, a quasi-uniform one.
\begin{enumerate}
\item[(a)] $s=2$: every prime $p$, with $\Lev=2\N$; the prime powers $q=8$, $9$, $16$, $25$, $27$, $32$, $49$, $64$, $81$, $121$, $125$, $128$, with $\Lev=2\N$.
\item[(b)] $s=3$: $q=3$, $q=4$, $q=7$ and every prime $p\ge11$, with $\Lev=3\N$; every prime $p\ge5$, with $\Lev=9\N$; the prime powers $q=25$, $32$, $49$, $64$, $81$, $121$, $125$, $128$, $169$, with $\Lev=3\N$. In particular every prime base $p\ge3$ is covered; $p=5$ only with $\Lev=9\N$.
\item[(c)] $s=4$: $q=4$, $49$, $121$, $125$ and every prime $p\ge17$, with $\Lev=4\N$; every prime $p\ge7$, with $\Lev=16\N$.
\item[(d)] $s=5$: $q=121$, $169$ and every prime $p\ge37$, with $\Lev=5\N$; every prime $p\ge11$, with $\Lev=25\N$.
\item[(e)] $s=6$: every prime $p\ge17$, with $\Lev=36\N$; $s=7$: $p=17,19,23$ and every prime $p\ge31$, with $\Lev=49\N$.
\item[(f)] every $s\ge2$: every prime $p>2^s$, with $\Lev=s\N$ and, by a second sequence, with $\Lev=s\N-1$; both are closed forms with written proofs.
\end{enumerate}
\end{theorem}

\begin{proof}
Each case is a cited result; Table~\ref{tab:coverage} lists the construction, the bound and the kind of proof for each.
(f) Corollary~\ref{cor:dyadic} (staggered residue sequence, $\Lev=s\N$) and Corollary~\ref{cor:korobov} with Corollary~\ref{cor:residueQU} (residue sequence, $\Lev=s\N-1$), both for every prime $p>2^s$.
(a) $p=2$: Theorem~\ref{thm:B2}; $p=3$: Remark~\ref{rem:B3} (computer-assisted), and Theorem~\ref{thm:block}(a) with $\Lev=4\N$ (proof); $p\ge5$: Theorem~\ref{thm:WF}. The prime powers: Proposition~\ref{prop:primepower}.
(b) $q=3$: Theorem~\ref{thm:S3}; $q=4$: Theorem~\ref{thm:K4}; $q=7$: Theorem~\ref{thm:p7}; $p\ge11$: case (f), since $11>2^3$. $\Lev=9\N$: Theorem~\ref{thm:block}(b), for every prime $p\ge5$. The prime powers: Proposition~\ref{prop:primepower}.
(c) $q=4$: Theorem~\ref{thm:S4}; $q=49,121,125$: Proposition~\ref{prop:primepower}; $p\ge17$: case (f). $\Lev=16\N$: Theorem~\ref{thm:block}(c) (written proof for $p\ge29$, exact evaluation for $7\le p\le109$).
(d) $q=121,169$: Proposition~\ref{prop:primepower}; $p\ge37$: case (f), since $37$ is the least prime above $2^5$. $\Lev=25\N$: Theorem~\ref{thm:block}(d), computer-assisted (the written proof of (c) covers $p\ge84$).
(e) Theorem~\ref{thm:block}(d), computer-assisted.
\end{proof}

\begin{table}[t]
\centering
\footnotesize
\caption{Coverage by dimension (Theorem~\ref{thm:coverage}): base, construction, level family $\Lev$ (the exponents $m$ with a certified bound on $\sep{P_m}$), the bound, and the kind of proof (``c.a.'': computer-assisted, with the scope stated in the cited result). The bound for every $N$ is $\LB\,q^{-\Delta_{\Lev}/s}$ with $\Delta_{\Lev}=s$ for $\Lev=s\N$ and $s\N-1$ and $\Delta_{\Lev}=s^2$ for $\Lev=s^2\N$ (Lemma~\ref{lem:prefix}). $\rho=\lfloor p^{1/s}\rfloor$. For $s=4,5$ the written proof of Theorem~\ref{thm:block}(c) gives $1-6/p$ and $1-17/p$ from $p=29$ and $84$ on. Constructions certified at every full level with constants tending to $1$ from large primes on (staggered windows, Section~\ref{sec:antipode}) are not repeated here.}
\label{tab:coverage}
\setlength{\tabcolsep}{3pt}
\begin{tabular}{@{}c l l l l l@{}}
\toprule
$s$ & base & construction & $\Lev$ & $\sep{P_m}\ge$ & proof\\
\midrule
2 & $q=2$ & Thm.~\ref{thm:B2} & $2\N$ & $0.3678$ & c.a.\\
 & $q=3$ & Rem.~\ref{rem:B3} & $2\N$ & $0.2669$ & c.a.\\
 & $p\ge5$ & brick wall, Thm.~\ref{thm:WF} & $2\N$ & $(p-3)/(2p)$ & proof\\
 & $p\ge5$ & residue seq., Cor.~\ref{cor:korobov} & $2\N-1$ & $(\rho-1)p^{-1/2}$ & proof\\
 & $p\ge3$ & block-periodic, Thm.~\ref{thm:block}(a) & $4\N$ & $1-2/p$ & proof\\
 & $q=8,\dots,128$ & Prop.~\ref{prop:primepower} & $2\N$ & Tab.~\ref{tab:primepower} & c.a.\\
\addlinespace
3 & $q=3$ & Thm.~\ref{thm:S3} & $3\N$ & $0.05598$ & c.a.\\
 & $q=4$ & Thm.~\ref{thm:K4} & $3\N$ & $0.11293$ & c.a.\\
 & $p=7$ & Thm.~\ref{thm:p7} & $3\N$ & $1/7$ & c.a.\\
 & $p\ge11$ & stagg.\ residue, Cor.~\ref{cor:dyadic} & $3\N$ & $(\lfloor p/4\rfloor-1)/p$ & proof\\
 & $p\ge11$ & residue seq., Cor.~\ref{cor:korobov} & $3\N-1$ & $(\rho-1)p^{-1/3}$ & proof\\
 & $p\ge5$ & block-periodic, Thm.~\ref{thm:block}(b) & $9\N$ & $1-2/p$ & proof\\
 & $q=25,\dots,169$ & Prop.~\ref{prop:primepower} & $3\N$ & Tab.~\ref{tab:primepower} & c.a.\\
\addlinespace
4 & $q=4$ & Thm.~\ref{thm:S4} & $4\N$ & $0.03013$ & c.a.\\
 & $p\ge17$ & stagg.\ residue, Cor.~\ref{cor:dyadic} & $4\N$ & $(\lfloor p/8\rfloor-1)/p$ & proof\\
 & $p\ge17$ & residue seq., Cor.~\ref{cor:korobov} & $4\N-1$ & $(\rho-1)p^{-1/4}$ & proof\\
 & $p\ge7$ & block-periodic, Thm.~\ref{thm:block} & $16\N$ & $1-4/p$ ($p\ge11$) & c.a.\\
 & $q=49,121,125$ & Prop.~\ref{prop:primepower} & $4\N$ & Tab.~\ref{tab:primepower} & c.a.\\
\addlinespace
5 & $p\ge37$ & stagg.\ residue, Cor.~\ref{cor:dyadic} & $5\N$ & $(\lfloor p/16\rfloor-1)/p$ & proof\\
 & $p\ge37$ & residue seq., Cor.~\ref{cor:korobov} & $5\N-1$ & $(\rho-1)p^{-1/5}$ & proof\\
 & $p\ge11$ & block-periodic, Thm.~\ref{thm:block} & $25\N$ & $1-8/p$ ($p\ge23$) & c.a.\\
 & $q=121,169$ & Prop.~\ref{prop:primepower} & $5\N$ & Tab.~\ref{tab:primepower} & c.a.\\
\addlinespace
6; 7 & see (e) & block-periodic, Thm.~\ref{thm:block} & $36\N$; $49\N$ & Tab.~\ref{tab:block} & c.a.\\
\addlinespace
$s$ & $p>2^s$ & stagg.\ residue, Cor.~\ref{cor:dyadic} & $s\N$ & $(\lfloor p/2^{s-1}\rfloor-1)/p$ & proof\\
 & $p>2^s$ & residue seq., Cor.~\ref{cor:korobov} & $s\N-1$ & $(\rho-1)p^{-1/s}$ & proof\\
 & $p\ge5a^*_s+4$ & block-periodic, Thm.~\ref{thm:block}(c) & $s^2\N$ & $1-(a^*_s+1)/p$ & proof ($s\ge4$)\\
\bottomrule
\end{tabular}
\end{table}

\begin{figure}[t]
\centering
\paperfig{fig_coverage}
\caption{Coverage for $s=2,\dots,7$ (Theorem~\ref{thm:coverage}). Each row is a level family: all full levels $s\N$, the prefixes $s\N-1$ one digit before them, and every $s$-th full level $s^2\N$ (which is contained in $s\N$, so every base covered in the first row is covered in the third). Dark: a written proof covers this base; light: computer-assisted only (the small-base theorems of this section, Theorem~\ref{thm:p7}, the exact evaluations of Theorem~\ref{thm:block}(c) below its proof range, and Theorem~\ref{thm:block}(d)); hatched: no $(0,s)$-sequence exists in base $q<s$. Every prime above $113$ is covered as in the last column, except for $s=7$, where the full levels are covered from $p=131$ on. Base $5$ in dimension three is certified only on $9\N$.}
\label{fig:dim3}
\end{figure}

\emph{What is not covered.} Bases $q<s$ are impossible: a $(0,s)$-sequence in base $q$ exists only for $s\le q$ \cite{Nie92}. Among the possible bases the following are not covered by this paper, and we do not claim that they cannot be: the prime powers that are not in Table~\ref{tab:primepower}, except $q=4$ for $s=3,4$ (for $q=8,9,16,27$ in dimension three the bound of Proposition~\ref{prop:primepower} is $0$, for every encoding resp.\ every $\F_p$-linear one); base $5$ in dimension three at every full level (Remark~\ref{rem:base5}); and for $s\ge4$ the primes $s\le p<2^s$ at every full level. On the levels $s^2\N$ these primes are covered from $p=7$, $11$, $17$ on for $s=4,5,6$, for $s=7$ at $p=17,19,23$ and $p\ge31$, and for $s\ge8$ where the proof range $p\ge5a^*_s+4$ of Theorem~\ref{thm:block}(c) starts below $2^s$ (not for $s=8,9,10$). Whether residue sequences are quasi-uniform for $s\le p<2^s$ is not addressed (Section~\ref{sec:open}).

\begin{remark}[Base five in dimension three]\label{rem:base5}
For $p=5$ and $s=3$, Theorem~\ref{thm:block}(b) gives $\sep{P_{9K}}\ge3/5$ for every $K$ (exact $\sep{P_9}=0.6211$), and hence $\delta(P_{[N]})\ge(3/5)\,5^{-3}N^{-1/3}$ for every $N$. A certificate at every full level is open. What was tried: (i) the ordered windows of an earlier version of this work (shown there to give no positive bound at $p=5$); (ii) the symmetric frames of this section have no three-dimensional instance over $\F_5$, since three places cannot form one orbit of a subgroup of the affine group of $\F_5$, whose order $20$ is prime to $3$; (iii) the staggered residue sequences have $\beta(k)\le0$ at $p=5$ (Corollary~\ref{cor:dyadic}), and among all their parameters up to the affine symmetries of $\F_5$ ($128$ choices of $\cstar$, $\theta$ and the weights) the best keeps $\sep{P_{3K}}\ge0.30$ for $K\le4$ (exact), but at deeper levels there are explicit close pairs: $\sep{P_{216}}\le0.0047$ for the choice that is best at $K\le4$, and $\sep{P_{300}}\le2.9\cdot10^{-6}$ for another (exact computations of the two points; whether any of them is quasi-uniform is open); (iv) for one-digit certificates that cover every possible cross response (the graph rule of Appendix~\ref{app:p7} and window coverings), and for several classes of two-digit certificates (ordered product windows as in (i) on $192$ relaxed instances, two further window classes on single instances, a frame-independent two-digit graph rule), a SAT solver reported that no positive bound exists in the instances solved; no independently checked certificate of these results was produced, so they are evidence about these classes only. General two-digit windows, certificates with more digits, and other constructions remain unexplored.
\end{remark}
